% Created 2026-03-17 ter 14:41
% Intended LaTeX compiler: lualatex
\documentclass[12pt]{article}
\usepackage[a4paper,left=3cm,top=3cm,right=2cm,bottom=2cm]{geometry}
\setcounter{secnumdepth}{4}
\setcounter{tocdepth}{2}
\usepackage{graphicx}
\usepackage{longtable}
\usepackage{wrapfig}
\usepackage{rotating}
\usepackage[normalem]{ulem}
\usepackage{amsmath}
\usepackage{amssymb}
\usepackage{capt-of}
\usepackage{hyperref}
\usepackage{fgv-header}
\author{Gustavo M. Mendes de Tarso}
\date{\today}
\title{}
\hypersetup{
 pdfauthor={Gustavo M. Mendes de Tarso},
 pdftitle={},
 pdfkeywords={},
 pdfsubject={},
 pdfcreator={Emacs 28.2 (Org mode 9.5.5)}, 
 pdflang={English}}
\begin{document}

\begingroup
\linespread{1}\selectfont
\begin{tcolorbox}[title=Ficha Técnica,
  colback=gray!5,colframe=gray!40,boxrule=0.4pt,sharp corners]
\begin{tblr}{rowsep=1pt,stretch=1, rows={t}, colspec={Q[l,2.8cm] X[l]}}
\textbf{Curso}                   & Mestrado de Políticas Públicas e Governo \\
\textbf{Turma}                   & T-01 \\
\textbf{Pólo}                    & Brasília \\
\textbf{Disciplina}              & Teorias de Administração Pública \\
\textbf{Professor}               & Bernardo Buta \\
\textbf{Trabalho}                & Aplicação do protocolo PRISMA 2020 \\
\textbf{Aluno(s)}                & Gustavo M. Mendes de Tarso \\
\textbf{Título do trabalho}      & Busca, seleção e síntese de artigos acadêmicos sobre teorias da administração pública \\
\end{tblr}
\end{tcolorbox}
\endgroup

\FGVNewPage

\section{Introdução}
\label{sec:orgb8eb422}

Este trabalho apresenta, em formato alinhado ao \textbf{PRISMA 2020}, o fluxo de identificação, triagem, elegibilidade e inclusão de artigos acadêmicos relacionados ao tema da disciplina \textbf{Teorias de Administração Pública}. Dada a amplitude do tema, a busca foi estruturada a partir dos eixos mais recorrentes na bibliografia da disciplina e no debate contemporâneo da área: \textbf{modelo burocrático}, \textbf{neo-institucionalismo}, \textbf{Nova Gestão Pública (NPM) e pós-NPM}, \textbf{governança pública} e \textbf{burocracia de nível de rua}.

A estratégia de identificação dos estudos combinou duas rotas complementares. A primeira consistiu na leitura e classificação da \textbf{bibliografia-base da disciplina}, já disponibilizada para a aula. A segunda consistiu em \textbf{busca acadêmica complementar} em páginas de periódicos, registros DOI, repositórios institucionais e metadados acadêmicos acessíveis, com o objetivo de localizar revisões, textos de síntese e artigos teóricos capazes de atualizar o debate.

Foram identificados \textbf{17 registros} ao todo. Após a remoção de \textbf{1 duplicata} e de \textbf{1 item não elegível por se tratar de livro, e não de artigo de periódico}, restaram \textbf{15 registros} para triagem. Em seguida, \textbf{3 registros} foram excluídos por baixa aderência ao recorte ou por excessiva especificidade temática. Os \textbf{12 textos remanescentes} foram avaliados em maior profundidade, dos quais \textbf{2} foram excluídos após leitura por apresentarem contribuição predominantemente normativa ou escopo demasiadamente estreito. Ao final, \textbf{10 artigos} foram incluídos na síntese qualitativa.

Convém registrar que esta aplicação do PRISMA 2020 tem finalidade \textbf{didática e metodológica de transparência}, isto é, busca relatar de modo claro como os estudos foram encontrados e selecionados, sem pretensão de exaustividade absoluta típica de uma revisão sistemática em larga escala.

\section{Estratégia de busca}
\label{sec:org9723b02}

\subsection{Pergunta orientadora}
\label{sec:orga4c6beb}

\begin{quote}
Quais artigos acadêmicos melhor representam e atualizam o debate teórico contemporâneo em administração pública nos eixos institucionalismo, burocracia, Nova Gestão Pública/pós-NPM, governança pública e implementação?
\end{quote}

\subsection{Expressões de busca utilizadas}
\label{sec:orgb585f27}

\begin{quote}
("public administration" OR "administração pública")
AND
("neo-institutionalism" OR "institutional theory" OR bureaucracy OR "street-level bureaucracy" OR "new public management" OR "post-new public management" OR "public governance" OR "governança pública")
AND
(article OR review OR "literature review" OR "systematic literature review")
\end{quote}

Como a busca foi realizada em fontes heterogêneas, e não em uma única base indexada, a expressão acima foi adaptada conforme o ambiente de busca, preservando sempre o mesmo núcleo terminológico.

\subsection{Fontes consultadas}
\label{sec:orgc0375a6}
\begin{itemize}
\item Bibliografia-base da disciplina
\item Páginas de periódicos científicos
\item Registros DOI
\item Repositórios institucionais
\item Metadados acadêmicos acessíveis durante a busca
\end{itemize}

\subsection{Critérios de inclusão}
\label{sec:org663fa94}
\begin{itemize}
\item Artigos de periódico
\item Aderência clara a pelo menos um dos eixos centrais da disciplina
\item Relevância teórica, de síntese ou de atualização bibliográfica
\item Disponibilidade de texto completo ou de metadados suficientes para avaliação inicial
\end{itemize}

\subsection{Critérios de exclusão}
\label{sec:org4961d98}
\begin{itemize}
\item Livros e capítulos
\item Duplicatas
\item Estudos excessivamente setoriais, sem contribuição substantiva ao debate teórico geral
\item Textos de natureza predominantemente normativa, ensaística ou programática, quando pouco úteis à síntese do campo
\end{itemize}

\section{PRISMA 2020}
\label{sec:org0c2ba6e}

\subsection{Fluxo textual dos estudos}
\label{sec:orga173ff2}

\subsubsection{Identificação}
\label{sec:org1884ef5}
\begin{itemize}
\item Registros identificados: \textbf{n = 17}
\begin{itemize}
\item Bibliografia-base da disciplina: \textbf{n = 9}
\item Busca acadêmica complementar: \textbf{n = 8}
\end{itemize}

\item Registros removidos antes da triagem: \textbf{n = 2}
\begin{itemize}
\item Duplicatas removidas: \textbf{n = 1}
\item Item removido por não ser artigo de periódico: \textbf{n = 1}
\end{itemize}
\end{itemize}

\subsubsection{Triagem}
\label{sec:org55c0b72}
\begin{itemize}
\item Registros triados: \textbf{n = 15}
\item Registros excluídos: \textbf{n = 3}
\begin{itemize}
\item Foco empírico excessivamente setorial: \textbf{n = 1}
\item Recorte excessivamente específico dentro do tema de governança: \textbf{n = 2}
\end{itemize}
\end{itemize}

\subsubsection{Elegibilidade}
\label{sec:orga1a103a}
\begin{itemize}
\item Relatórios buscados para recuperação: \textbf{n = 12}
\item Relatórios não recuperados: \textbf{n = 0}

\item Relatórios avaliados para elegibilidade: \textbf{n = 12}
\item Relatórios excluídos após leitura em texto completo: \textbf{n = 2}
\begin{itemize}
\item Contribuição mais normativa do que teórica: \textbf{n = 1}
\item Escopo estreito para o objetivo da síntese: \textbf{n = 1}
\end{itemize}
\end{itemize}

\subsubsection{Inclusão}
\label{sec:org235300f}
\begin{itemize}
\item Estudos incluídos na síntese qualitativa: \textbf{n = 10}
\end{itemize}

\subsection{Tabela-resumo do fluxo PRISMA 2020}
\label{sec:org4fa26b1}

\begin{table}[htbp]
\centering
\scriptsize
\setlength{\tabcolsep}{4pt}
\renewcommand{\arraystretch}{1.15}
\begin{tabularx}{\textwidth}{>{\raggedright\arraybackslash}X >{\raggedleft\arraybackslash}m{0.18\textwidth}}
\toprule
Etapa PRISMA 2020 & n \\
\midrule
Registros identificados & 17 \\
Registros removidos antes da triagem & 2 \\
Registros triados & 15 \\
Registros excluídos & 3 \\
Relatórios buscados para recuperação & 12 \\
Relatórios não recuperados & 0 \\
Relatórios avaliados para elegibilidade & 12 \\
Relatórios excluídos após leitura completa & 2 \\
Estudos incluídos na síntese qualitativa & 10 \\
\bottomrule
\end{tabularx}
\normalsize
\end{table}

\subsection{Estudos classificados no fluxo}
\label{sec:orgca26e86}

A tabela abaixo considera os \textbf{15 registros efetivamente triados}. A duplicata removida e o livro excluído antes da triagem não aparecem nesta classificação.

\begin{table}[htbp]
\centering
\scriptsize
\setlength{\tabcolsep}{3pt}
\renewcommand{\arraystretch}{1.15}
\begin{tabularx}{\textwidth}{>{\raggedright\arraybackslash}X >{\raggedright\arraybackslash}m{0.16\textwidth} >{\raggedright\arraybackslash}m{0.18\textwidth} >{\raggedright\arraybackslash}m{0.26\textwidth}}
\toprule
Estudo & Origem & Situação no fluxo & Motivo \\
\midrule
Institutionalized Organizations: Formal Structure as Myth and Ceremony & Bibliografia-base / busca complementar & Incluído & Clássico da teoria institucional organizacional \\
The Iron Cage Revisited: Institutional Isomorphism and Collective Rationality in Organizational Fields & Bibliografia-base & Incluído & Texto central para isomorfismo institucional e campo organizacional \\
Political Science and the Three New Institutionalisms & Bibliografia-base & Incluído & Síntese fundamental das três vertentes do neo-institucionalismo \\
A Public Management for All Seasons? & Bibliografia-base & Incluído & Marco teórico da Nova Gestão Pública \\
Public Governance: Balancing Stakeholder Power in a Network Society & Busca complementar & Incluído & Referência central sobre governança pública \\
Modelos organizacionais e reformas da administração pública & Bibliografia-base & Incluído & Articulação entre burocracia, gerencialismo e governança no debate brasileiro \\
The manifold meanings of post-New Public Management & Bibliografia-base & Incluído & Revisão de literatura sobre os sentidos do pós-NPM \\
Trends in Public Administration after Hegemony of the New Public Management & Busca complementar & Incluído & Atualização do debate sobre pós-NPM e governança \\
Street-Level Bureaucracy in Public Administration: A systematic literature review & Busca complementar & Incluído & Revisão contemporânea sobre burocracia de nível de rua \\
The relevance of neo-institutionalism for organizational change & Busca complementar & Incluído & Atualiza o institucionalismo para análise de mudança organizacional \\
Programa Cultura Viva e o campo organizacional da cultura & Bibliografia-base & Excluído na triagem & Foco empírico setorial, com menor centralidade para o mapa teórico geral da disciplina \\
New paths for public governance: a systematic literature review and integrative framework & Busca complementar & Excluído na triagem & Recorte excessivamente específico dentro do eixo de governança \\
Strategic public value(s) governance: a systematic literature review and future research agenda & Busca complementar & Excluído na triagem & Foco estreito em subcampo específico de governança pública \\
Reforma da Nova Gestão Pública: agora na agenda da América Latina, no entanto... & Bibliografia-base & Excluído na elegibilidade & Texto historicamente relevante, porém mais ensaístico-normativo que sintético para o recorte \\
Developing New Public Governance as a public management reform model & Busca complementar & Excluído na elegibilidade & Contribuição predominantemente programática e normativa \\
\bottomrule
\end{tabularx}
\normalsize
\end{table}

\subsection{Representação esquemática em PlantUML}
\label{sec:org8a41c68}

\begin{center}
\includesvg[width=.9\linewidth]{prisma_teorias_adm_publica_fluxo}
\end{center}

\section{Síntese dos artigos selecionados}
\label{sec:org0e12734}

Os \textbf{10 artigos incluídos} permitem organizar o campo em três grandes blocos analíticos. O primeiro bloco é o do \textbf{neo-institucionalismo}. Nele, \emph{Institutionalized Organizations: Formal Structure as Myth and Ceremony} e \emph{The Iron Cage Revisited} cumprem papel fundacional ao explicar por que organizações adotam estruturas formais, rotinas e símbolos não apenas por eficiência instrumental, mas também por busca de \textbf{legitimidade}. \emph{Political Science and the Three New Institutionalisms} amplia esse quadro ao distinguir institucionalismo histórico, da escolha racional e sociológico, oferecendo base conceitual importante para compreender a pluralidade do institucionalismo. Já \emph{The relevance of neo-institutionalism for organizational change} atualiza essa tradição ao mostrar como o institucionalismo pode ser mobilizado para analisar transformação organizacional, e não apenas estabilidade.

O segundo bloco reúne os textos sobre \textbf{burocracia, Nova Gestão Pública, pós-NPM e governança pública}. \emph{A Public Management for All Seasons?} permanece como referência para a compreensão da NPM, sobretudo por sistematizar valores como eficiência, competição, mensuração de resultados e descentralização gerencial. \emph{Modelos organizacionais e reformas da administração pública} é especialmente relevante para o contexto da disciplina por articular, em uma mesma chave analítica, o modelo burocrático, o gerencialismo e a governança pública, mostrando continuidades e tensões entre eles. \emph{The manifold meanings of post-New Public Management} e \emph{Trends in Public Administration after Hegemony of the New Public Management} demonstram que o debate pós-NPM não substituiu inteiramente a NPM, mas passou a reordenar o campo com maior ênfase em coordenação, integração, colaboração e arranjos híbridos. \emph{Public Governance: Balancing Stakeholder Power in a Network Society} reforça essa inflexão ao situar a governança pública como abordagem centrada em redes, pluralismo de atores, coordenação horizontal e redistribuição de poder entre Estado e sociedade.

O terceiro bloco é o da \textbf{implementação e burocracia de nível de rua}. Nesse eixo, \emph{Street-Level Bureaucracy in Public Administration: A systematic literature review} mostra a permanência da agenda inaugurada por Lipsky, agora revisitada à luz da produção recente. O artigo evidencia que a burocracia de nível de rua segue central para entender como políticas públicas são efetivamente concretizadas, pois a implementação depende da discricionariedade, das rotinas e das condições de trabalho dos agentes que interagem diretamente com os cidadãos.

Tomados em conjunto, os artigos selecionados sugerem que a administração pública continua sendo um campo fortemente \textbf{interdisciplinar}, no qual teorias oriundas da sociologia, da ciência política, dos estudos organizacionais e da gestão pública são constantemente mobilizadas para interpretar Estado, organizações públicas, reformas, implementação e governança. Ao mesmo tempo, a seleção mostra que o debate contemporâneo não se resume à oposição entre burocracia e gerencialismo. O que emerge com maior força é um campo em que \textbf{instituições, legitimidade, coordenação, discricionariedade, redes e governança} se entrelaçam na análise da ação pública.

\section{Considerações finais}
\label{sec:org986e458}

À luz do fluxo PRISMA 2020 aqui aplicado, verifica-se que o tema da disciplina é amplo e exige uma estratégia de busca que combine \textbf{clássicos fundacionais} com \textbf{artigos contemporâneos de revisão e atualização}. O resultado final confirma que uma seleção equilibrada de estudos não deve privilegiar apenas textos históricos, nem apenas revisões recentes, mas articular ambos para reconstruir o desenvolvimento do campo.

A busca também mostrou que parte significativa da literatura útil à disciplina está distribuída entre tradições analíticas distintas. O neo-institucionalismo oferece instrumentos poderosos para compreender estabilidade, mudança e legitimidade organizacional. A literatura sobre NPM e pós-NPM ajuda a interpretar reformas, desempenho e gestão. A governança pública amplia a lente para arranjos colaborativos, pluralismo e redes. Por sua vez, a burocracia de nível de rua recoloca a implementação no centro da análise, lembrando que políticas públicas não existem apenas no plano formal, mas se realizam concretamente na interação entre burocratas e cidadãos.

Em síntese, o corpus selecionado é adequado para sustentar discussões em sala de aula e para servir de base a trabalhos acadêmicos introdutórios na disciplina \textbf{Teorias de Administração Pública}, pois oferece simultaneamente \textbf{fundamentos conceituais}, \textbf{mapa das principais correntes} e \textbf{atualizações recentes do campo}.
\end{document}